\section{Related Work}

\subsection{Token Log-Probability Baselines}

Token-level log-probabilities are among the most accessible uncertainty signals in autoregressive LLMs, requiring no additional inference beyond the standard forward pass. \citet{fadeeva2024} introduce the Claim Conditioned Probability (CCP), which uses mean log-probability while conditioning on the factual claim portion of a response. CCP achieves AUROC 0.72--0.80 on TriviaQA and Natural Questions across seven LLMs and four languages, establishing mean log-prob as a strong implicit baseline for factual hallucination detection. \citet{kadavath2022} demonstrate that token log-probabilities correlate with model self-knowledge via the P(True) and P(IK) constructs, laying the empirical groundwork for log-probability uncertainty signals. Our work extends this line by showing that CCP's implicit choice of \emph{mean} aggregation is suboptimal for factual recall benchmarks by a margin of 5--12 AUROC points when min is used instead --- a gap that CCP could have closed without architectural modification.

\subsection{Multi-Sample Consistency Methods}

\citet{manakul2023} introduce SelfCheckGPT, which detects hallucination by measuring semantic consistency across $N$ sampled outputs ($N{=}20$ for competitive performance). Variants using NLI, BERTScore, or N-gram overlap achieve AUROC 0.72--0.78 on sentence-level hallucination detection. SelfCheckGPT includes ad hoc token log-probability baselines without systematic ablation of aggregation function or directional predictions. \citet{farquhar2024} introduce Semantic Entropy, which clusters $N$ sampled outputs by semantic equivalence (via NLI) and computes predictive entropy over clusters; it achieves AUROC $\approx$0.79 on TriviaQA with $N{=}10$ samples. SE's use of sum (unnormalized within-cluster entropy) is an implicit aggregation choice that is not ablated. A critical practical contrast: our best single-pass result (raw\_sum AUROC 0.895--0.896 on TriviaQA, under our evaluation protocol) exceeds published SE AUROC at one-tenth the inference cost, highlighting the unexplored potential of zero-cost single-pass aggregations.

Recent extensions reduce SE's sampling overhead: \citet{ciosek2025} achieve matching AUROC with 53\% fewer samples via Bayesian SE, and \citet{nguyen2025} improve via pairwise semantic similarity (SNNE). These methods remain in the multi-sample regime and inherit the implicit aggregation choice. Our work is orthogonal: we address the single-pass regime and provide the first principled aggregation selection criterion within it.

\subsection{Supervised and Probe-Based Methods}

\citet{shelmanov2025} train UQ heads on attention maps to detect hallucination at claim level across Mistral, LLaMA, and Gemma architectures, achieving state-of-the-art performance at the cost of supervised training per model. \citet{moslonka2025} propose EPR, which extracts features from top-$k$ logits and applies a learned detector; it achieves strong performance on QA but requires training data. \citet{ma2025} and \citet{zhang2025} address related challenges in logit-level uncertainty and robust uncertainty quantification. Our work is entirely unsupervised: no fine-tuning, no learned components, no labeled calibration data. \textit{Note: \citet{ma2025}, \citet{moslonka2025}, and \citet{zhang2025} are arXiv preprints; citations are included for completeness and should be verified against published versions prior to submission.}

\subsection{Hallucination Benchmarks and Taxonomy}

\citet{lin2022} introduce TruthfulQA, a benchmark of 817 questions designed to elicit imitative-falsehood hallucinations. \citet{farquhar2024} provide reproducible binary-label splits of TriviaQA and Natural Questions from the semantic\_uncertainty repository. Our taxonomy --- hallucination type $\to$ distribution shape $\to$ aggregation optimality --- is the first mechanistic classification that predicts optimal aggregation function from hallucination type alone.

\subsection{Aggregation Function Comparison: The Open Gap}

\citet{liu2025} survey uncertainty quantification and calibration for LLMs, identifying systematic comparison of token-level log-probability aggregation strategies as an open challenge. Our paper directly fills this gap: controlled ablation with mechanism-grounded directional predictions confirmed with bootstrap CIs across two model families.
